source cutoff.


\section{Corrected audit constants and the closed-sector ledger}\label{sec:auditconstants}

The independent hostile audit accepted the rough-modulus proof and its physical
transference, but identified three corrections that must be incorporated into
any authoritative manuscript.  A fourth endpoint split was requested for full
reviewability.  None changes an exponent or reopens a closure node.

\begin{center}
\small
\begin{tabularx}{\textwidth}{@{}p{0.20\textwidth}X X@{}}
\toprule
Issue & Correct statement & Consequence \\
\midrule
Principal coprimality removal
& Use full inclusion--exclusion:
  $O_j(2^j(D/P+1))$, not the literal coefficient $j(D/P+1)$.
& For the physical value $j=5$, only the absolute constant changes. \\
\addlinespace
Odd-M\"obius recursion
& The dyadic identity \cref{eq:odd-recursion} is charged one additional
  factor $O(\log D)$ in the conservative global ledger.
& Replace $\LL^{113+3b/2}$ by $\LL^{114+3b/2}$. \\
\addlinespace
Cellwise frequency legality
& Define the cutoff with
  $Q_C^{\min}=\min_{(A,k)\in C}Q_{A,k,C}$ and impose
  $2H_C<Q_C^{\min}$.
& Every retained frequency is legal on every line in the cell. \\
\addlinespace
Complementary-modulus endpoint
& Split $fe\le R$, $R<fe<2R$, and $fe\ge2R$ before using
  $\#\{r:e\mid r,\,R\le fr<2R\}\le2R/(fe)$.
& Makes \cref{eq:r-count} explicit; no extra term or logarithm appears. \\
\bottomrule
\end{tabularx}
\end{center}

The kernel-checked outer CARD5 coefficient and the crude source bounds give the
fixed coefficient
\begin{equation}\label{eq:fixed-coeff}
                  252\cdot5!\cdot6!=21{,}772{,}800.
\end{equation}
The inherited source compiler contributes $\LL^{111}$; one dyadic partition in
$k$, the explicit $\log k$, and the conservative odd-M\"obius recursion each
cost at most one further logarithm.  If $H\le\LL^b$, the frequency summation
costs $\LL^{3b/2}$.  A sufficient global envelope for the newly closed region
is therefore
\begin{equation}\label{eq:corrected-ledger}
 \begin{split}
 C\,21{,}772{,}800\,X\LL^{114+3b/2}
 \bigl(&X^{-\delta+o(1)}+\LL^{-J}\\
       &+P_*^{-1}+D^{-1}\bigr).
 \end{split}
\end{equation}
For a requested output $X\LL^{-B}$, choose, for example,
\begin{equation}\label{eq:J-choice}
                    J>B+115+\frac{3b}{2}.
\end{equation}
The strict extra unit is a harmless safety margin.  All terms other than
$\LL^{-J}$ are fixed-power small on the closed region and absorb the complete
finite logarithmic ledger.

\classification{\newapplication}{Independent audit correction; conditional on the inherited $\LL^{111}$ source ledger.}
\begin{proposition}[Corrected arbitrary-log closure ledger]\label{prop:corrected-ledger}
Under \cref{hyp:frozen}, \cref{eq:corrected-ledger} is a sufficient uniform
envelope for every fixed-polylogarithmic band closed by
\cref{thm:transference-main}.  In particular, the three audit corrections do
not invalidate the closure of the former Sector~A.
\end{proposition}

\begin{proof}
Insert \cref{eq:delta} into \cref{eq:transference-main}, multiply by the
fixed coefficient \cref{eq:fixed-coeff} and by the four declared logarithmic
sources, and use $H^{3/2}\le\LL^{3b/2}$.  The choice
\cref{eq:J-choice} leaves at least $\LL^{-B-1}$ after the conservative ledger;
the fixed-power terms dominate every fixed power of $\LL$ for sufficiently
large $X$.
\end{proof}

\section{Adaptive conductor cutoff}\label{sec:adaptive}

The old constants $\eta=\rho=1/10000$ were useful for an initial partition but
are not intrinsic to the physical source.  The variance theorem supplies the
cellwise squared nonprincipal deficit
\begin{equation}\label{eq:LambdaC}
 \Lambda_C=
 \left(1+\frac{K_C}{R_C}\right)
 \left(\frac{R_C^2}{X}+\frac{R_C}{K_CP_C}\right).
\end{equation}
For $K_C\le R_C$, the leading multiplicity factor is bounded by $2$; retaining
it uniformly avoids a second definition at $K_C=R_C$.

\classification{\newreduction}{Independently audited source-specific definition; the minimum-conductor correction is authoritative.}
\begin{definition}[Cellwise legal cutoff]\label{def:HC}
Fix a requested output exponent $S>0$, a fixed logarithmic reserve $c_0>0$,
and a fixed cap $b>0$.  Put
\begin{equation}\label{eq:Qmin}
            Q_C^{\min}=\min_{(A,k)\in C}Q_{A,k,C}.
\end{equation}
If there is a positive integer $H$ satisfying
\begin{equation}\label{eq:HC-conditions}
 H\le\LL^b,
 \qquad 2H<Q_C^{\min},
 \qquad
 (1+H)^{3/2}\sqrt{\Lambda_C}\le\LL^{-(S+c_0)},
\end{equation}
let $H_C$ be the largest such $H$.  If there is no such positive integer, set
$H_C=0$.
\end{definition}

The definition is admissible rather than globally optimal.  It adapts to the
actual conductor deficit and never asks a line to support a Fourier frequency
outside its cyclic Nyquist range.  The reserve $c_0$ is chosen to dominate the
finite coefficient and logarithmic compiler costs; it is not a new analytic
parameter.

If $R_C^2/X$ is already fixed-power small and $K_C\le R_C$, the full cap
$H_C=\lfloor\LL^b\rfloor$ is admitted by the sufficient inequality
\begin{equation}\label{eq:full-cap-condition}
          \frac{K_C}{R_C}
          \gg P_C^{-1}\LL^{2(S+c_0)+3b}.
\end{equation}
Thus the controlled low-frequency region can extend to an inverse physical
prime factor in $K_C/R_C$, rather than stopping at an arbitrary inverse
logarithm.

\classification{\newapplication}{Independent analytic audit; conditional on the finite compiler ledger.}
\begin{proposition}[Low-pass estimate at the adaptive cutoff]\label{prop:adaptive-low}
For $H_C$ as in \cref{def:HC}, the complete physical first-moment contribution
of the frequencies $0<|\xi|\le H_C$, summed over all cells, is
\begin{equation}\label{eq:adaptive-low-bound}
                        O_S(X\LL^{-S}).
\end{equation}
The principal term is included.
\end{proposition}

\begin{proof}
Apply \cref{thm:transference-main} cell by cell.  The nonprincipal term is at
most $\LL^{-(S+c_0)}$ times the natural cell mass by
\cref{eq:HC-conditions}.  Choose $c_0$ larger than the declared finite
compiler cost and sum the finitely many dyadic/Mellin cells.  For the principal
term, choose $J$ as in \cref{eq:J-choice}; the $P_C^{-1}$ and $D_C^{-1}$ terms
are fixed-power small in the physical range.  If $H_C=0$, the low-pass
contribution is empty.
\end{proof}

\section{Exact high-pass reassembly}\label{sec:reassembly}

For each line, let $P_C^{\mathrm{lo}}$ denote the finite Fourier projection
\begin{equation}\label{eq:Plo}
  [P_C^{\mathrm{lo}}b^\perp](t)
   =\sum_{0<|\xi|_Q\le H_C}c_{A,k,C}(\xi)\eord{Q}{\xi t},
\end{equation}
where $|\xi|_Q$ is the least absolute representative.  The zero mode is absent
because $b^\perp$ is centred.  The high-pass target is
$(I-P_C^{\mathrm{lo}})b^\perp$.

The finite coefficient $c_{C,A,k}$ below retains the five-prime source weight,
all signs, the smoothing and Mellin labels, and every local source restriction.
It is not a generic bounded coefficient.

\classification{\newreduction}{Exact linear decomposition; independently audited physical transference; conditional on the frozen full-source compiler.}
\begin{theorem}[Source-exact high-pass reassembly]\label{thm:reassembly}
Define
\begin{equation}\label{eq:Rhi}
 \boxed{
 \Rhi(X)=252\sum_{C,A,k}c_{C,A,k}\log k
   \sum_{t\in\cI_{A,k,C}}\mu(d_t)
      [(I-P_C^{\mathrm{lo}})b^\perp_{A,k,C}](t).
 }
\end{equation}
Then, for every fixed requested $S$,
\begin{equation}\label{eq:R11-Rhi}
 \boxed{
          \Roneone(X)=\Rhi(X)+O_S(X\LL^{-S})
 }
\end{equation}
under \cref{hyp:frozen}.
\end{theorem}

\begin{proof}
On every line,
\[
 b^\perp=P_C^{\mathrm{lo}}b^\perp
            +(I-P_C^{\mathrm{lo}})b^\perp
\]
is a literal finite Fourier identity.  Substitute it into the frozen
first-moment $R_{11}$ row before taking any absolute value or square.  The sum
of the low-pass terms is $O_S(X\LL^{-S})$ by
\cref{prop:adaptive-low}; all previously owned source pieces are already
included in the same error by \cref{hyp:frozen}.  What remains is exactly
\cref{eq:Rhi}.  No coefficient, ordering state, or local label is removed.
\end{proof}

\begin{warning}[Mixed-frequency firewall]\label{fire:mixed}
If a two-copy argument is applied after the split, then
\[
 |L^{\mathrm{lo}}+L^{\mathrm{hi}}|^2
 =|L^{\mathrm{lo}}|^2+|L^{\mathrm{hi}}|^2
 +2\Re(L^{\mathrm{lo}}\overline{L^{\mathrm{hi}}}).
\]
The proof of \cref{thm:reassembly} does not claim that the cross term vanishes.
It avoids the issue by proving the complete low contribution small at
first-moment level and replacing the original row by the exact high-pass row.
A future MHHC estimate must start from that high-pass row.
\end{warning}

The old three-sector formulation
\[
 \cO_{\mathrm{bal,low}}^{\mathrm{upgrade}}
 +\cO_{\mathrm{bal,high}}+\cO_{\mathrm{unbal}}
\]
is therefore superseded.  It was a useful diagnostic split, but the controlling
frontier is now the single signed functional \cref{eq:Rhi}.

\section{The remaining $\Rhi$ problem and quantified nonclosures}\label{sec:Rhi}

The exact open estimate is
\begin{equation}\label{eq:Rhi-target}
             \boxed{\Rhi(X)\ll_S X\LL^{-S}}
\end{equation}
for one sufficiently large fixed $S$ required by the downstream finite ledger.
In balanced cells, $\Rhi$ contains only frequencies above the chosen fixed
polylogarithmic cap.  In strongly unbalanced cells, \cref{def:HC} may give
$H_C=0$, in which case every nonzero frequency remains.

The following calculations explain why several natural architectures do not
yet prove \cref{eq:Rhi-target}.  They are capacity statements about particular
proofs, not lower bounds for $\Rhi$ and not impossibility theorems.

\subsection{Common-frequency cardinal interpolation}

For one line define the exact band-limited interpolation
\begin{equation}\label{eq:cardinal-g}
 g_i(x)=\sum_{t\in\cI_i}b_i^\perp(t)
  \operatorname{sinc}\left(\frac{Dx-a_i}{q}-t\right),
 \qquad \operatorname{sinc}(u)=\frac{\sin\pi u}{\pi u}.
\end{equation}
At $x=(a_i+qt)/D$, this equals $b_i^\perp(t)$.  Its Fourier transform is
\begin{equation}\label{eq:cardinal-G}
 G_i(\tau)=\frac qD\ee{-a_i\tau/D}
   \1_{|\tau|\le D/(2q)}
   \sum_tb_i^\perp(t)\ee{-q\tau t/D},
\end{equation}
and Plancherel gives
\begin{equation}\label{eq:cardinal-energy}
      \int_{\R}|G_i(\tau)|^2\,d\tau
       =\frac qD\sum_t|b_i^\perp(t)|^2\ll B_0^2.
\end{equation}
For each common $\tau$, \cref{thm:roughvariance} applies to
$a_n=\mu(n)\ee{\tau n/D}$ without a derivative loss.  On $|\tau|\le H$ this
yields
\begin{equation}\label{eq:cardinal-bound}
 \sum_i\left|\int_{|\tau|\le H}G_i(\tau)E_i(\tau)\,d\tau\right|
 \ll_jB_0\sqrt{RK}\sqrt{H\left(DR+D^2/P\right)}.
\end{equation}
At the full bandwidth $H\asymp T=D/R$, normalisation by $DK$ gives
